% ============================================================
%  Interpreting FCFS Appointment Simulation Results Under
%  the Current Formulation
%  Metric behavior, access losses, and model checks
% ============================================================

\documentclass[11pt]{article}

\usepackage[margin=1.0in]{geometry}
\usepackage{amsmath}
\usepackage{amssymb}
\usepackage{booktabs}
\usepackage{graphicx}
\usepackage{float}
\usepackage{hyperref}
\usepackage{array}
\usepackage[table]{xcolor}
\usepackage{microtype}
\usepackage{parskip}
\usepackage{enumitem}
\usepackage{caption}

\hypersetup{colorlinks=true,linkcolor=black,urlcolor=blue,citecolor=black}

\graphicspath{%
  {../../figures/}%
  {../../../../../outputs/class1_balking/figures/}%
  {../../../../../outputs/class1_balking_threshold/figures/}%
}

\sloppy
\setcounter{tocdepth}{2}

\newcommand{\mainfigure}[3]{%
  \begin{figure}[H]
    \centering
    \includegraphics[width=\linewidth]{#1}
    \caption{#2}
    \label{#3}
  \end{figure}%
}
\newcommand{\code}[1]{\texttt{#1}}

\title{\textbf{Interpreting FCFS Appointment Simulation Results\\
       Under the Current Formulation}\\[4pt]
       \large Metric behavior, access losses, and model checks}
\author{CUIMC Appointment Simulation}
\date{May 2026}

\begin{document}
\maketitle
\tableofcontents
\newpage

\section{Abstract}

This note interprets simulation outputs from a first-contact FCFS
appointment model.  It is an \emph{interpretation note}, not an
empirical claim about real clinic operations.  All statements are
conditional on the current model formulation, and numbers should be
read as model outputs, not clinical measurements.

\medskip

\textbf{Main finding.}  Under the current formulation, average
utilization and the overall served rate separate sharply.  At
baseline, average utilization is \textbf{0.839} while the overall
served rate is \textbf{0.269}.  This gap is mechanically plausible
rather than paradoxical: utilization is normalized by the number of
available slots, while the served rate is normalized by the number of
patient arrivals.  With 100 arrivals per day and 32 slots per day the
system is structurally oversubscribed.  Even if every scheduled slot
were completed without any no-show, only $32/100 = 0.32$ of arrivals
could be served.  The observed served rate of 0.269 is further reduced
below this ceiling by no-shows.

\medskip

\textbf{Baseline loss decomposition.}
Of all arrivals at baseline, approximately 0.356 balk (reject the
offered appointment), 0.323 cancel after booking, and 0.052 are
no-shows.  No-offer (horizon-full rejection) is zero at baseline,
confirming that the booking horizon is not saturated under baseline
demand.  The low served rate is therefore driven primarily by balking
and post-booking cancellation, not by horizon saturation.

\medskip

\textbf{Offered vs.\ accepted wait.}
Mean offered wait (9.30 days at baseline) includes both patients who
accepted and patients who later balked.  Mean accepted wait (8.35
days) excludes balked patients and is therefore a selected,
downward-biased measure of the delay experience.  Offered wait should
always be reported alongside the no-offer share, because patients who
receive no offer are excluded from both wait metrics.

\section{Simulation Mechanics Relevant for Interpretation}

The points below are necessary for reading the metric outputs
correctly.  They describe the current implementation rather than
design aspirations.

\begin{itemize}[leftmargin=1.5em, itemsep=3pt]
  \item \textbf{Finite booking horizon.}  The simulation maintains a
    rolling calendar of $H = 14$ residual days.  When all days are
    fully booked, the next arrival receives no offer.
  \item \textbf{Two patient classes.}  Each class has its own Poisson
    arrival rate, balking rule, no-show rule, and cancellation probability.
  \item \textbf{Class-specific Poisson arrivals.}
    $M_1^D \sim \mathrm{Poisson}(\lambda_1)$ and
    $M_2^D \sim \mathrm{Poisson}(\lambda_2)$ independently each day.
  \item \textbf{Arrivals pooled and randomly permuted.}  All arrivals
    are combined and randomly shuffled into a single FCFS queue.  No
    class is given scheduling priority.
  \item \textbf{Earliest available slot offered.}  Each patient is
    offered the smallest $r \in \{1,\ldots,H-1\}$ with remaining
    capacity.  The offered delay $\tau$ equals this residual day $r$.
  \item \textbf{Offered delay $\tau$ recorded before balking.}
    The value $\tau$ is logged at the time of the offer, regardless of
    whether the patient subsequently accepts or balks.
  \item \textbf{Balking depends on $\tau$.}  A class-$i$ patient
    accepts with probability $1 - b_i(\tau)$, where $b_i(\cdot)$ is a
    step function with higher rejection above a class-specific threshold.
    A balked patient leaves permanently.
  \item \textbf{No-show depends on original $\tau$.}  No-show
    probability $\xi_i(\tau)$ is anchored to the \emph{original}
    offered delay, not the remaining wait on the day of service.
  \item \textbf{Future cancellations only.}  At the start of each day,
    each booked patient on a future residual day ($r \geq 1$) cancels
    independently with probability $\phi_i$.  Canceled slots are
    returned to the pool and can be rebooked.
  \item \textbf{No same-day cancellations.}  Patients at $r = 0$
    either attend or are no-shows.
  \item \textbf{No rebooking of no-show slots.}  No-show slots count
    as lost capacity.
  \item \textbf{No-offer rule when horizon is full.}  A patient with
    no slot available is recorded as \code{no\_offer} and excluded
    from the offered-wait distribution.
  \item \textbf{Utilization counts completed visits per slot.}
    $\frac{1}{D}\sum_d \frac{\text{served slots}_d}{S}$.
    No-shows do not contribute.
\end{itemize}

\section{Baseline Diagnosis}

\subsection{Baseline Parameter Setup}

\begin{itemize}[leftmargin=1.5em, itemsep=2pt]
  \item Two symmetric classes: $\lambda_1 = \lambda_2 = 50$ per day,
    total $\lambda = 100$.
  \item $S = 32$ slots per day; $H = 14$-day booking horizon.
  \item Cancellation probability: $\phi_1 = \phi_2 = 0.10$.
  \item Balking threshold: 9 days; high-delay balking probability:
    0.50; low: 0.00 (step = 0.50).
  \item No-show threshold: 6 days; high-delay no-show probability:
    0.30; low: 0.00 (step = 0.30).
\end{itemize}

\subsection{Baseline Metric Table}

\begin{table}[H]
\centering
\renewcommand{\arraystretch}{1.18}
\begin{tabular}{lrr}
\toprule
Metric & Baseline & Scenario~2 \\
\midrule
Average utilization & 0.839 & 1.000 \\
Overall served rate & 0.269 & 0.395 \\
Mean accepted wait (days) & 8.35 & 4.29 \\
Mean offered wait (days) & 9.30 & 5.06 \\
Class served-rate gap ($\Delta_{\text{access}}$) & 0.001 & $-$0.008 \\
Class delay gap ($\Delta_{\text{delay}}$) & 0.003 & 0.005 \\
\bottomrule
\end{tabular}
\caption{Baseline and Scenario~2 summary metrics.  Scenario~2 changes
  several assumptions simultaneously and is a comparison point, not a
  one-parameter causal test.}
\end{table}

\subsection{Utilization--Access Gap}

\[
  \text{average utilization} = \frac{\text{completed visits}}{\text{available slots}},
  \qquad
  \text{overall served rate} = \frac{\text{completed visits}}{\text{arrivals}}.
\]

\subsection{Demand-to-Capacity Calculation}

With $\lambda = 100$ and $S = 32$, the demand-to-capacity ratio is
$100/32 \approx 3.1$.  Even with zero no-shows and zero cancellations,
the served rate is bounded by $32/100 = 0.32$.  The observed 0.269
falls below this ceiling because no-shows eliminate completed visits
after scheduling.

\section{Arrival Outcome Decomposition}

\begin{table}[H]
\centering
\renewcommand{\arraystretch}{1.18}
\begin{tabular}{lrl}
\toprule
Outcome & Share & Mechanism \\
\midrule
Served   & 0.269 & Accepted, did not cancel, attended \\
Balked   & 0.356 & Offered a slot; rejected due to long delay \\
Canceled & 0.323 & Accepted; canceled before service day \\
No-show  & 0.052 & Accepted; not canceled; did not attend \\
No-offer & 0.000 & Horizon full; no slot offered \\
\midrule
Total & $\approx 1.000$ & \\
\bottomrule
\end{tabular}
\caption{Arrival outcome decomposition at baseline.}
\end{table}

\textbf{Baseline is not horizon-saturated.}
No-offer is zero, so access failure is not caused by the horizon filling.

\textbf{Low served rate is driven by balking and cancellation.}
Together they account for $0.356 + 0.323 = 0.679$ of arrivals.
Reducing no-shows alone would not materially raise the served rate.

\textbf{High cancellation is mechanically plausible.}
A patient booked 10 days ahead faces about 10 independent cancellation
draws.  With $\phi = 0.10$, the survival probability is
$(0.90)^{10} \approx 0.35$.  This should be verified in code to
confirm cancellation is applied \emph{once per residual day per
patient}, not once per booking.

\mainfigure{fcfs_arrival_outcome_decomposition.png}{%
  Arrival outcome decomposition under demand pressure.
  At baseline ($\lambda = 100$), no-offer is zero.
  No-offer becomes substantial only above approximately $\lambda = 110$.
  Balking and cancellation dominate at baseline.}{fig:outcome-decomp}

\section{Cross-Metric Diagnostics}

\mainfigure{metric_access_drivers.png}{%
  Overall served rate as demand changes.  Class~1 varies on the x-axis
  while Class~2 stays fixed at baseline.}{fig:access-drivers}

\mainfigure{fcfs_capacity_stress_curves.png}{%
  FCFS capacity stress curves ($\lambda = 30$ to $170$, 50/50 mix,
  three seeds averaged).  Utilization stays near 0.84 across a wide
  demand range while served rate falls steeply.  No-offer becomes
  material only above approximately $\lambda = 110$.}{fig:stress-curves}

Three diagnostic patterns emerge.

\textbf{1.~Utilization can stay near 0.84 while served rate falls.}
Slots are still occupied by patients who accept; completed visits per
slot stays high even as the served fraction of arrivals collapses.

\textbf{2.~Accepted wait is selected downward by balking.}
Patients who would have waited longest have left the system.  Offered
wait includes those rejected patients and is the less biased measure.

\textbf{3.~No-offer becomes important only above $\lambda \approx 110$.}
At baseline demand no-offer is zero.  The no-offer share must be
reported separately whenever it is non-negligible.

\section{Metric Hierarchy Under This Formulation}

\begin{table}[H]
\centering
\renewcommand{\arraystretch}{1.25}
\small
\begin{tabular}{p{0.16\textwidth}p{0.20\textwidth}p{0.26\textwidth}p{0.27\textwidth}}
\toprule
Metric & What it measures & Interpretation & Main caveat \\
\midrule
Overall served rate
  & Fraction of arrivals completing a visit
  & Primary access metric
  & Bounded above by $S/\lambda$ \\[4pt]
Average utilization
  & Fraction of slots filled by completed visits
  & Capacity-use metric
  & Not an access metric \\[4pt]
Mean offered wait
  & Avg delay offered to all patients receiving an offer
  & Broader congestion measure
  & Excludes no-offer patients \\[4pt]
Mean accepted wait
  & Avg delay among patients who booked
  & Patient-level delay for those who accepted
  & Selected downward by balking \\[4pt]
Balking rate
  & Fraction of offered patients rejecting the delay
  & Diagnostic for rejected offers
  & Not a success metric \\[4pt]
Class served-rate gap
  & Class~1 served rate minus Class~2
  & Class-level access disparities
  & Near zero under symmetric parameters \\
\bottomrule
\end{tabular}
\caption{Metric hierarchy under the current FCFS formulation.}
\end{table}

\section{Sensitivity Screen: Main Patterns}

\mainfigure{metric_utilization_drivers.png}{%
  Average utilization driver plots.  No-show behavior is the clearest
  utilization driver.}{fig:util-drivers}

\mainfigure{metric_wait_drivers.png}{%
  Mean offered wait driver plots.  Arrival pressure is the strongest
  driver.  Higher balking or cancellation can shorten offered wait for
  reasons that worsen access rather than improve it.}{fig:wait-drivers}

\textbf{Demand pressure} raises offered wait and reduces served rate
monotonically.  Utilization is relatively insensitive across a wide
demand range.

\textbf{No-show behavior} is the primary utilization driver.
The regression screen (Appendix~\ref{app:regression}) identifies the
no-show threshold as the strongest single predictor of utilization.

\textbf{Cancellation} has ambiguous aggregate effects.  Cancelled
slots are returned to the pool and can be rebooked, so higher
cancellation can shorten offered wait while still reducing completed
visits per original patient.  The regression screen shows a
\emph{positive} coefficient for cancellation on both utilization and
served rate, which is counterintuitive and should be verified at code
level (Section~\ref{sec:validation}).

\textbf{Balking} can reduce accepted wait while reducing access.
This is the main argument for reporting offered wait rather than
accepted wait as the primary delay metric.

\textbf{Class gaps} arise from asymmetric behavioral parameters.
Cancellation probability differences are the largest class-access
driver in the regression screen.

\section{Balking Deep Dive: Selection vs.\ Access}

\mainfigure{percent_serviced_by_class.png}{%
  Class~1 high-balking-probability sweep: served rate by class.
  As Class~1 rejection probability rises, Class~1 access falls while
  Class~2 access rises slightly.}{fig:balk-deep-served}

\mainfigure{mean_accepted_delay_by_class.png}{%
  Class~1 high-balking-probability sweep: mean accepted delay by class.
  Class~1 accepted delay falls as balking rises because long-delay
  patients exit the accepted sample.  This is a selection effect.}{fig:balk-deep-accepted}

\textbf{Main point.}  Increasing balking can make accepted wait appear
to improve while access worsens.  Offered wait and served rate together
are necessary to interpret the change correctly.

\section{Class-Gap Interpretation Under FCFS}

Under the current formulation, class labels are \emph{exchangeable}
when parameters are symmetric.  Arrivals are pooled and randomly
permuted; no class receives FCFS priority.

\textbf{Symmetric parameters imply expected class gap near zero.}
The observed baseline gap of 0.001 is consistent with simulation noise.

\textbf{Gap drivers} in order of absolute standardized coefficient
(Appendix~\ref{app:regression}):
\begin{enumerate}[leftmargin=2em, itemsep=3pt]
  \item Cancellation probability gap: $\hat{\beta} = -0.452$
  \item Balking threshold gap: $\hat{\beta} = +0.429$
  \item Balking step gap: $\hat{\beta} = -0.349$
  \item No-show threshold gap: $\hat{\beta} = +0.268$
  \item No-show step gap: $\hat{\beta} = -0.211$
  \item Class~1 arrival share: $\hat{\beta} = -0.120$
\end{enumerate}

\section{What Seems Mechanically Robust vs.\ What Needs Validation}
\label{sec:validation}

\subsection{Mechanically Robust}

\begin{itemize}[leftmargin=1.5em, itemsep=4pt]
  \item Utilization and served rate have different denominators.
    High utilization can coexist with low served rate by definition.
  \item The demand ceiling: served rate cannot exceed $32/100 = 0.32$
    without overbooking.
  \item Offered wait is less selected than accepted wait.
  \item No-offer is a distinct access failure excluded from wait metrics.
  \item Class labels are exchangeable under symmetric assumptions.
  \item No-offer is zero at baseline; failures are behavioral, not
    capacity.
\end{itemize}

\subsection{Needs Validation or Calibration}

\begin{itemize}[leftmargin=1.5em, itemsep=4pt]
  \item \textbf{Cancellation applied once per residual day.}
    Verify in code: one draw per patient per day, not one draw per
    booking.
  \item \textbf{Positive cancellation coefficient on served rate.}
    Counterintuitive; verify what fraction of canceled slots are
    rebooked.
  \item \textbf{No-offer propagation within a day.}
    Confirm all remaining queue patients receive no-offer when horizon
    is full; confirm permutation is re-drawn each day.
  \item \textbf{No-show slots never rebooked.}  Deliberate choice;
    should be stated explicitly when comparing to literature models.
  \item \textbf{Same-day cancellations excluded.}  All same-day losses
    route through no-show.
  \item \textbf{Class-gap stability with more seeds.}  Current grids
    use one seed per grid point.
  \item \textbf{Realism of the demand/capacity ratio.}  3.1 arrivals
    per slot is a stress setting, not a calibrated clinic baseline.
  \item \textbf{Behavioral parameter calibration.}  Balking step 0.50,
    no-show step 0.30, and cancellation 0.10 are assumed, not
    estimated.
\end{itemize}

\section{Summary for Discussion}

The simulation output suggests the following under the current
formulation.

Utilization near 0.84 and served rate near 0.27 can coexist because
they measure different things.  The structural ratio
$\lambda/S \approx 3.1$ constrains served rate below 0.32 even before
behavioral losses.  Balking and cancellation reduce it further and
dominate no-show in the baseline decomposition.

The metric ordering is: (1) overall served rate for access, (2) average
utilization for slot use, (3) mean offered wait for patient-facing
delay, (4) no-offer share as a separate indicator, (5) mean accepted
wait as a secondary selected metric, (6) balking rate and class gap as
diagnostics.

The sensitivity screen is mechanically consistent with this ordering.
Several findings should be checked before outputs support design
decisions.

\bigskip
\textbf{Questions for discussion:}
\begin{enumerate}[leftmargin=2em, itemsep=4pt]
  \item Is the baseline demand/capacity ratio intentional as a stress
    scenario, or should it be recalibrated?
  \item Is the positive cancellation coefficient verified at code level?
    What fraction of canceled slots are rebooked?
  \item Should same-day cancellations be allowed?
  \item Should no-show slots be rebooked?
  \item Are threshold-jump step functions the right behavioral model,
    or should smoother functions be explored?
  \item Are the class behavioral parameters population-specific or
    sensitivity levers only?
\end{enumerate}

\appendix

\section{Exact Metric Definitions and Notation}
\label{app:metrics}

Let $i \in \{1,2\}$, $D$ = measured days, $S$ = slots/day, $H$ =
horizon length. For class $i$: $A_i$ arrivals, $B_i$ booked,
$K_i$ balked, $O_i = B_i + K_i$ offered, $N_i$ no-offer,
$C_i$ canceled, $H_i$ no-shows, $V_i$ served.
$A_i = V_i + C_i + H_i + K_i + N_i$.

\begin{align*}
  \code{percent\_serviced}_i &= V_i / A_i, \\
  \code{overall\_percent\_serviced} &= \textstyle\sum_i V_i / \sum_i A_i, \\
  \code{average\_utilization} &= \frac{1}{D}\sum_d \text{served slots}_d / S, \\
  \code{mean\_offered\_booking\_delay} &= \textstyle\sum_i \text{total offered delay}_i / \sum_i O_i, \\
  \code{mean\_accepted\_booking\_delay} &= \textstyle\sum_i \text{total accepted delay}_i / \sum_i B_i, \\
  \code{balking\_rate} &= \textstyle\sum_i K_i / \sum_i O_i, \\
  \Delta_{\text{access}} &= \code{percent\_serviced}_1 - \code{percent\_serviced}_2, \\
  \Delta_{\text{delay}} &= \bar{\tau}^{\text{offered}}_2 - \bar{\tau}^{\text{offered}}_1, \\
  \Delta_{\text{balking}} &= K_1/O_1 - K_2/O_2.
\end{align*}

\textbf{Behavioral step functions:}
\[
  b_i(\tau) = \begin{cases}
    b_i^{\text{low}} & \tau < \theta_i^{\text{balk}}, \\
    b_i^{\text{high}} & \tau \geq \theta_i^{\text{balk}}.
  \end{cases}
\]
Balking step $= b_i^{\text{high}} - b_i^{\text{low}}$.
No-show $\xi_i(\tau)$ has the same structure.

\section{Sensitivity Design}
\label{app:sensitivity-design}

All grids centered on \code{configs/baseline.yaml}.
One-axis slices vary Class~1 while Class~2 stays at baseline.
Two-class heatmaps vary both classes independently.
Each grid point uses seed 2027; multi-seed comparisons average
seeds 2027, 2028, 2029.

\begin{table}[H]
\centering
\renewcommand{\arraystretch}{1.15}
\begin{tabular}{p{0.26\textwidth}rrr}
\toprule
Parameter & Baseline & Range & Grid size \\
\midrule
Arrivals (per class), 1D slices & 50 & 25--75 & continuous \\
Total arrivals, FCFS stress & 100 & 30--170 & continuous \\
Balking step & 0.50 & 0.00--1.00 & 21 \\
Balking threshold (days) & 9 & 0--13 & 14 \\
No-show step & 0.30 & 0.00--1.00 & 21 \\
No-show threshold (days) & 6 & 0--13 & 14 \\
Cancellation probability & 0.10 & 0.00--0.30 & 21 \\
\bottomrule
\end{tabular}
\caption{Sensitivity parameter ranges.}
\end{table}

\textbf{Regression screen:} 240 randomized settings (seed 9137);
two seeds (4101, 4102) averaged per setting; 12 features (total
arrival rate, Class~1 share, average level and Class~1$-$Class~2
gap for each behavioral parameter); OLS with HC3 robust SEs;
only $p < 0.05$ coefficients shown.

\textbf{Balking deep dives:} Class~1 high-balking-probability sweep
over $\{0.0, 0.1, \ldots, 0.9\}$; Class~1 threshold sweep over
$\{3,\ldots,13\}$; 100 seeds per point; Class~2 fixed at baseline.

\section{Supplementary Sensitivity Plots}
\label{app:sensitivity-plots}

\mainfigure{no_show_step_interaction_heatmap.png}{No-show step heatmap. Left: utilization. Right: Class~1 served-rate gap.}{fig:app-noshow-step}
\mainfigure{no_show_threshold_interaction_heatmap.png}{No-show threshold heatmap. Left: utilization. Right: Class~1 gap.}{fig:app-noshow-thresh}
\mainfigure{cancellation_probability_interaction_heatmap.png}{Cancellation heatmap. Left: served rate. Right: Class~1 gap.}{fig:app-cancel}
\mainfigure{balking_step_interaction_heatmap.png}{Balking step heatmap. Left: served rate. Right: Class~1 gap.}{fig:app-balk-step}
\mainfigure{balking_threshold_interaction_heatmap.png}{Balking threshold heatmap. Left: offered wait. Right: Class~1 gap.}{fig:app-balk-thresh}
\mainfigure{arrival_mix_interaction_heatmap.png}{Arrival mix heatmap. Left: served rate. Right: offered wait.}{fig:app-arrival-mix}
\mainfigure{scenario_metric_comparison.png}{Scenario-level metric comparison (baseline vs.\ Scenario~2).}{fig:app-scenario}
\mainfigure{balking_threshold_jump_heatmaps.png}{Common balking threshold and jump level.}{fig:app-balk-common}
\mainfigure{no_show_threshold_jump_heatmaps.png}{Common no-show threshold and jump level.}{fig:app-noshow-common}

\section{Regression Screen}
\label{app:regression}

\begin{table}[H]
\centering
\small
\renewcommand{\arraystretch}{1.08}
\begin{tabular}{p{0.17\textwidth}p{0.27\textwidth}rp{0.20\textwidth}r}
\toprule
Target & Feature & Std.\ coef. & 95\% HC3 CI & HC3 $p$ \\
\midrule
Utilization & Avg no-show threshold & $+0.512$ & $[0.433,\;0.591]$ & $<0.001$ \\
Utilization & Avg no-show step & $-0.423$ & $[-0.501,\;-0.346]$ & $<0.001$ \\
Utilization & Avg cancellation prob. & $+0.320$ & $[0.247,\;0.393]$ & $<0.001$ \\
Utilization & Total arrival rate & $-0.109$ & $[-0.188,\;-0.031]$ & $0.006$ \\
Utilization & Avg balking threshold & $-0.086$ & $[-0.162,\;-0.009]$ & $0.028$ \\
\midrule
Overall served rate & Total arrival rate & $-0.774$ & $[-0.855,\;-0.694]$ & $<0.001$ \\
Overall served rate & Avg no-show threshold & $+0.251$ & $[0.193,\;0.308]$ & $<0.001$ \\
Overall served rate & Avg no-show step & $-0.175$ & $[-0.232,\;-0.119]$ & $<0.001$ \\
Overall served rate & Avg cancellation prob. & $+0.117$ & $[0.058,\;0.176]$ & $<0.001$ \\
\midrule
Offered wait & Total arrival rate & $+0.576$ & $[0.499,\;0.653]$ & $<0.001$ \\
Offered wait & Avg cancellation prob. & $-0.441$ & $[-0.506,\;-0.376]$ & $<0.001$ \\
Offered wait & Avg balking threshold & $+0.278$ & $[0.202,\;0.353]$ & $<0.001$ \\
Offered wait & Avg balking step & $-0.206$ & $[-0.291,\;-0.122]$ & $<0.001$ \\
\midrule
Class gap & Cancellation prob.\ gap & $-0.452$ & $[-0.558,\;-0.347]$ & $<0.001$ \\
Class gap & Balking threshold gap & $+0.429$ & $[0.331,\;0.528]$ & $<0.001$ \\
Class gap & Balking step gap & $-0.349$ & $[-0.431,\;-0.266]$ & $<0.001$ \\
Class gap & No-show threshold gap & $+0.268$ & $[0.179,\;0.357]$ & $<0.001$ \\
Class gap & No-show step gap & $-0.211$ & $[-0.286,\;-0.136]$ & $<0.001$ \\
Class gap & Class~1 arrival share & $-0.120$ & $[-0.209,\;-0.030]$ & $0.009$ \\
\bottomrule
\end{tabular}
\caption{Significant standardized regression coefficients ($p<0.05$, HC3 robust SEs).
  Gap = Class~1 minus Class~2 value.}
\end{table}

\mainfigure{regression_significant_standardized_coefficients.png}{%
  Significant standardized regression coefficients with 95\% HC3
  confidence intervals.}{fig:app-regression}

\end{document}
